% Options for packages loaded elsewhere
\PassOptionsToPackage{unicode}{hyperref}
\PassOptionsToPackage{hyphens}{url}
\documentclass[
  11pt,
]{article}
\usepackage{xcolor}
\usepackage[left=3cm,right=3cm,top=2.5cm,bottom=2.5cm]{geometry}
\usepackage{amsmath,amssymb}
\setcounter{secnumdepth}{5}
\usepackage{iftex}
\ifPDFTeX
  \usepackage[T1]{fontenc}
  \usepackage[utf8]{inputenc}
  \usepackage{textcomp} % provide euro and other symbols
\else % if luatex or xetex
  \usepackage{unicode-math} % this also loads fontspec
  \defaultfontfeatures{Scale=MatchLowercase}
  \defaultfontfeatures[\rmfamily]{Ligatures=TeX,Scale=1}
\fi
\usepackage{lmodern}
\ifPDFTeX\else
  % xetex/luatex font selection
\fi
% Use upquote if available, for straight quotes in verbatim environments
\IfFileExists{upquote.sty}{\usepackage{upquote}}{}
\IfFileExists{microtype.sty}{% use microtype if available
  \usepackage[]{microtype}
  \UseMicrotypeSet[protrusion]{basicmath} % disable protrusion for tt fonts
}{}
\usepackage{setspace}
\makeatletter
\@ifundefined{KOMAClassName}{% if non-KOMA class
  \IfFileExists{parskip.sty}{%
    \usepackage{parskip}
  }{% else
    \setlength{\parindent}{0pt}
    \setlength{\parskip}{6pt plus 2pt minus 1pt}}
}{% if KOMA class
  \KOMAoptions{parskip=half}}
\makeatother
\usepackage{longtable,booktabs,array}
\usepackage{caption}
\captionsetup[table]{skip=6pt}
\newcounter{none} % for unnumbered tables
\usepackage{calc} % for calculating minipage widths
% Correct order of tables after \paragraph or \subparagraph
\usepackage{etoolbox}
\makeatletter
\patchcmd\longtable{\par}{\if@noskipsec\mbox{}\fi\par}{}{}
\makeatother
% Allow footnotes in longtable head/foot
\IfFileExists{footnotehyper.sty}{\usepackage{footnotehyper}}{\usepackage{footnote}}
\makesavenoteenv{longtable}
\usepackage{graphicx}
\makeatletter
\newsavebox\pandoc@box
\newcommand*\pandocbounded[1]{% scales image to fit in text height/width
  \sbox\pandoc@box{#1}%
  \Gscale@div\@tempa{\textheight}{\dimexpr\ht\pandoc@box+\dp\pandoc@box\relax}%
  \Gscale@div\@tempb{\linewidth}{\wd\pandoc@box}%
  \ifdim\@tempb\p@<\@tempa\p@\let\@tempa\@tempb\fi% select the smaller of both
  \ifdim\@tempa\p@<\p@\scalebox{\@tempa}{\usebox\pandoc@box}%
  \else\usebox{\pandoc@box}%
  \fi%
}
% Set default figure placement to htbp
\def\fps@figure{htbp}
\makeatother
\setlength{\emergencystretch}{3em} % prevent overfull lines
\providecommand{\tightlist}{%
  \setlength{\itemsep}{0pt}\setlength{\parskip}{0pt}}
%% tools/stamp.tex
%% Provenance footer for PDFs rendered from this project.
%% Vendored by zzcollab; this file is not project-specific. The
%% footer prints the source document and the time of rendering.
%% The git version is defined here too, and so is available to a
%% project that wants it on the page, but it is left off the
%% default footer: it already names the staged copy in share/ and
%% appears in its manifest row. All values are supplied at render
%% time by tools/stamp-render.R, which writes a small generated
%% file that redefines the macros below. The \providecommand
%% defaults keep a bare LaTeX compile working when the document is
%% built without the wrapper.
\usepackage{fancyhdr}
\usepackage{xcolor}
\providecommand{\stampsource}{\detokenize{(source not recorded)}}
\providecommand{\stamptime}{(time not recorded)}
\providecommand{\stampversion}{\detokenize{(version not recorded)}}
\setlength{\footskip}{40pt}
\newcommand{\stampfootcontent}{%
  \footnotesize\color{gray}%
  \parbox[b]{\textwidth}{%
    \stamptime\hfill\thepage\\[2pt]%
    \ttfamily\raggedright\stampsource}}
\pagestyle{fancy}
\fancyhf{}
\fancyfoot[C]{\stampfootcontent}
\renewcommand{\headrulewidth}{0pt}
\renewcommand{\footrulewidth}{0.4pt}
%% The title page uses the `plain` page style; redefine it so the
%% stamp appears there too.
\fancypagestyle{plain}{%
  \fancyhf{}%
  \fancyfoot[C]{\stampfootcontent}%
  \renewcommand{\headrulewidth}{0pt}%
  \renewcommand{\footrulewidth}{0.4pt}}
\renewcommand{\stampsource}{\textasciitilde{}/\allowbreak{}prj/\allowbreak{}res/\allowbreak{}04-runin-power-analysis/\allowbreak{}runinpower/\allowbreak{}analysis/\allowbreak{}report/\allowbreak{}04-gls-vs-ttest/\allowbreak{}bullets.Rmd}
\renewcommand{\stamptime}{2026-08-16 16:15 PDT}
\renewcommand{\stampversion}{\detokenize{276ac7f-wip}}
\usepackage{amsmath}
\usepackage{amssymb}
\usepackage{booktabs}
\usepackage{longtable}
\usepackage{float}
\usepackage[mathlines]{lineno}
\linenumbers
\DeclareMathOperator{\Var}{Var}
\DeclareMathOperator{\Cov}{Cov}
\DeclareMathOperator{\E}{E}
\usepackage{bookmark}
\IfFileExists{xurl.sty}{\usepackage{xurl}}{} % add URL line breaks if available
\urlstyle{same}
% fallback for those not using the hyperref driver hyperxmp:
\makeatletter
\@ifundefined{xmpquote}{\newcommand{\xmpquote}[1]{#1}}{}
\makeatother
\hypersetup{
  pdftitle={Structured summary: \textbar{} GLS versus Averaged Change-Score Estimators for Clinical Trials with Run-In Observations},
  pdfauthor={Ronald G. Thomas, Ph.D.},
  hidelinks,
  pdfcreator={LaTeX via pandoc}}

\title{Structured summary: \textbar{} GLS versus Averaged Change-Score Estimators for Clinical Trials with Run-In Observations}
\author{Ronald G. Thomas, Ph.D.\thanks{Wertheim School of Public Health, UCSD; ORCID: 0000-0003-1686-4965}}
\date{August 16, 2026}

\begin{document}
\maketitle

{
\setcounter{tocdepth}{2}
\tableofcontents
}
\setstretch{1.1}
\emph{This document is a structured, section-by-section summary of the key points argued in the companion manuscript, ``| GLS versus Averaged Change-Score Estimators for Clinical Trials with Run-In Observations.'' It is provided as a supplementary reading aid and does not stand on its own; all notation, derivations, simulation detail, and citations are in the main manuscript.}

\bigskip\noindent

\rule{\textwidth}{0.4pt}\bigskip

\section{Introduction}\label{introduction}

\textbf{The optimality of GLS is contingent on knowing the covariance, which in practice must be estimated and may be unstable in small samples.}

\begin{itemize}\setlength{\itemsep}{2pt}\setlength{\parskip}{0pt}
\item GLS is the best linear unbiased estimator (BLUE) only when the covariance matrix is known.
\item Feasible GLS substitutes estimated variance components, so its efficiency depends on estimation accuracy.
\item Small samples or short follow-up yield unstable estimates that can erode the theoretical advantage of GLS.
\end{itemize}

\textbf{The relative merits of change-score and ANCOVA analyses have been studied extensively across several methodological traditions.}

\begin{itemize}\setlength{\itemsep}{2pt}\setlength{\parskip}{0pt}
\item Formal efficiency comparisons for pretest-posttest designs trace back to @brogan1980 and @yang2001.
\item @frison1992 and @frison1997 established the efficiency ordering among post-treatment, change-score, and ANCOVA analyses with multiple baseline measurements, the closest antecedent of the present comparison.
\item Subsequent work contrasted baseline-as-covariate with baseline-as-response and developed semiparametric covariate adjustment.
\item Later contributions clarified when longitudinal models gain power over ANCOVA and adapted these ideas to run-in trials.
\end{itemize}

\textbf{The averaged change-score estimator is distinct from and less efficient than the constrained longitudinal data analysis estimator, which is excluded from the present comparison.}

\begin{itemize}\setlength{\itemsep}{2pt}\setlength{\parskip}{0pt}
\item The cLDA estimator places baseline in the response vector under a common-mean constraint.
\item It exploits the full covariance structure without specifying the random-effects distribution, ranking between the averaged estimator and GLS.
\item A full comparison including cLDA lies beyond the present scope but constitutes a natural extension.
\end{itemize}

\textbf{The longpower package implements the GLS baseline but omits the simpler estimators and the misspecification analysis developed here.}

\begin{itemize}\setlength{\itemsep}{2pt}\setlength{\parskip}{0pt}
\item It is the CRAN reference for slope-comparison power calculations under the Liu-Liang, Diggle, and Edland formulations.
\item Its routines correspond to the GLS estimator under correct specification, matching the Paper 1 baseline.
\item It does not cover ANCOVA, the averaged change-score test, or the efficiency loss from misspecified covariance.
\end{itemize}

\subsection{GLS under the mixed model}\label{gls-under-the-mixed-model}

\textbf{The GLS variance for the treatment effect follows from the relevant diagonal element of the inverse information matrix under the marginal covariance.}

\begin{itemize}\setlength{\itemsep}{2pt}\setlength{\parskip}{0pt}
\item The per-pair variance is the treatment-effect element of $(X'\Sigma^{-1}X)^{-1}$.
\item The marginal covariance decomposes as residual plus random-effects components.
\item The inverse is obtained efficiently via the Woodbury identity.
\end{itemize}

\subsection{ANCOVA with run-in mean as covariate}\label{ancova-with-run-in-mean-as-covariate}

\textbf{The ANCOVA variance is governed by the run-in to post-randomization correlation, which determines the magnitude of the variance reduction.}

\begin{itemize}\setlength{\itemsep}{2pt}\setlength{\parskip}{0pt}
\item The factor $(1 - \rho^2)$ scales the variance by the squared marginal correlation between the phase means.
\item With no run-in observations the covariate vanishes and ANCOVA reduces to the averaged estimator.
\item A large correlation, arising when the random-slope variance dominates the residual, yields substantial variance reduction.
\end{itemize}

\textbf{ANCOVA occupies a middle ground by estimating the baseline coefficient from the data while still using the predictive run-in observations.}

\begin{itemize}\setlength{\itemsep}{2pt}\setlength{\parskip}{0pt}
\item Unlike GLS, the coefficient on the run-in mean is estimated empirically rather than derived from assumed variance components.
\item Unlike the averaged estimator, it exploits the predictive value of the run-in observations.
\item It therefore reduces variance without committing to a full random-effects specification.
\end{itemize}

\subsection{BLUP as theoretical benchmark}\label{blup-as-theoretical-benchmark}

\textbf{The three estimators are linear functions of the change scores and admit a common benchmark in the best linear unbiased predictor of the random slopes.}

\begin{itemize}\setlength{\itemsep}{2pt}\setlength{\parskip}{0pt}
\item Each estimator can be viewed as a linear combination of the observed change scores.
\item The BLUP of the subject-specific random slope provides the natural reference predictor.
\item It weights the residuals by the inverse marginal covariance using the GLS fixed-effect estimate.
\end{itemize}

\textbf{GLS attains the minimum variance among linear unbiased estimators when the covariance is correct, while the simpler estimators are specific approximations to this benchmark.}

\begin{itemize}\setlength{\itemsep}{2pt}\setlength{\parskip}{0pt}
\item The group-specific slope estimators inherit the variance of the treatment-effect fixed effect.
\item GLS is optimal among linear unbiased estimators under a correctly specified marginal covariance.
\item The averaged and ANCOVA estimators are particular linear approximations to this optimal benchmark.
\end{itemize}

\textbf{Misspecification erodes but, in the settings studied here, does not eliminate the GLS efficiency advantage over the simpler approximations.}

\begin{itemize}\setlength{\itemsep}{2pt}\setlength{\parskip}{0pt}
\item GLS propagates errors in the assumed covariance through its inverse weighting.
\item The averaged and ANCOVA estimators depend on fewer assumed variance components.
\item Section 4 quantifies the resulting erosion; in the misspecifications examined, GLS retains a large margin.
\end{itemize}

\subsection{\texorpdfstring{Misspecified \(G\)}{Misspecified G}}\label{misspecified-g}

\textbf{Misestimating the random-slope variance forfeits the optimality of GLS while leaving the averaged estimator unchanged.}

\begin{itemize}\setlength{\itemsep}{2pt}\setlength{\parskip}{0pt}
\item An incorrect random-slope variance renders the assumed-covariance GLS no longer minimum-variance.
\item The error enters through the inverse-covariance weighting of GLS.
\item The averaged change-score estimator's definition does not use this variance component, so its variance does not depend on the assumed value.
\item In the sweep below, however, the misspecified GLS variance peaks at 5.30 while the averaged variance on the $\gamma$ scale is 8.89: the GLS margin is eroded, not eliminated.
\end{itemize}

\subsection{\texorpdfstring{Misspecified \(R\)}{Misspecified R}}\label{misspecified-r}

\textbf{Misspecifying the residual correlation biases the GLS variance machinery but again leaves the averaged estimator intact.}

\begin{itemize}\setlength{\itemsep}{2pt}\setlength{\parskip}{0pt}
\item Assuming compound symmetry when the truth is AR(1) is a representative misspecification.
\item Both the GLS variance formula and the Woodbury correction become biased.
\item The averaged estimator's definition does not depend on the assumed correlation, although its $\gamma$-scale variance still changes with the true $R$ (8.4 to 9.1 over the plotted range).
\item The misspecified GLS sandwich variance stays four- to elevenfold below the averaged variance across the sweep.
\end{itemize}

\subsection{Design}\label{design}

\textbf{A factorial Monte Carlo design complements the analytical results by varying sample size, design, signal-to-noise, and analysis method under a compound-symmetric truth.}

\begin{itemize}\setlength{\itemsep}{2pt}\setlength{\parskip}{0pt}
\item Sample size and the run-in design configuration are crossed.
\item The signal-to-noise ratio is varied; the true correlation is compound symmetric throughout.
\item Three analysis methods are compared: correctly specified GLS, misspecified GLS, and the averaged change-score test.
\end{itemize}

\textbf{Each configuration is replicated 1000 times and evaluated on empirical variance and power; coverage is not assessed.}

\begin{itemize}\setlength{\itemsep}{2pt}\setlength{\parskip}{0pt}
\item One thousand trials are simulated per configuration.
\item Empirical standard deviation and empirical power are recorded for each estimator.
\item Coverage of nominal 95\% confidence intervals, empirical size at $\Delta = 0$, larger sample sizes, higher ratios, an AR(1) truth arm, and 5000 replicates remain future work.
\end{itemize}

\subsection{Results}\label{results}

\textbf{The simulation confirms the analytic ordering and quantifies the cost of misspecifying the random-slope variance.}

\begin{itemize}\setlength{\itemsep}{2pt}\setlength{\parskip}{0pt}
\item The empirical standard deviation of the GLS estimator is the smallest of the three under correct specification.
\item On the $\gamma$ scale the averaged estimator is nearly as efficient as GLS in the no-run-in design, but carries two to three times the GLS standard deviation in the run-in designs at $\sigma_b/\sigma = 1$, growing to eight to fifteen times at $\sigma_b/\sigma = 5$.
\item Misspecifying the random-slope variance at most roughly doubles the GLS standard deviation, which remains well below that of the averaged estimator in every run-in cell.
\end{itemize}

\section{Discussion}\label{discussion}

\textbf{GLS is the efficient choice under a correctly specified, adequately estimated model, which in practice means moderate-to-large samples.}

\begin{itemize}\setlength{\itemsep}{2pt}\setlength{\parskip}{0pt}
\item GLS attains the minimum variance among linear unbiased estimators when the covariance is known.
\item Its advantage is realized only when the variance components are estimated accurately, which requires sufficient sample size and follow-up.
\item In small samples the instability of the estimated covariance can erase the theoretical gain.
\end{itemize}

\textbf{The misspecifications studied here erode the GLS advantage without overturning it; the open robustness question concerns feasible GLS with estimated components.}

\begin{itemize}\setlength{\itemsep}{2pt}\setlength{\parskip}{0pt}
\item GLS weights observations by the assumed inverse covariance, so errors in $G$ or $R$ propagate into its sampling variance.
\item In the sweeps of Section 4 the misspecified GLS sandwich variance peaks at 5.30 against a $\gamma$-scale averaged-estimator variance of 8.89, and in the simulation the misspecified GLS standard deviation remains well below the averaged standard deviation in every run-in cell.
\item Whether feasible GLS, with variance components estimated from small samples, retains this margin is not studied here.
\end{itemize}

\textbf{The efficiency gap between GLS and the simpler estimators narrows as the design grows, so the choice matters most in small, information-poor designs.}

\begin{itemize}\setlength{\itemsep}{2pt}\setlength{\parskip}{0pt}
\item With many run-in and treatment observations the phase means become near-sufficient for the random slope.
\item The averaged and ANCOVA estimators then recover most of the information GLS extracts, and the relative efficiency approaches one.
\item The estimator choice is therefore most consequential when observations are few and the slope-to-noise ratio is high.
\end{itemize}

\textbf{Two extensions, the seemingly unrelated regression estimator and the constrained longitudinal data analysis estimator, would complete the efficiency picture.}

\begin{itemize}\setlength{\itemsep}{2pt}\setlength{\parskip}{0pt}
\item The SUR formulation offers a middle ground between the averaged estimator and full GLS but adds implementation complexity.
\item The cLDA estimator exploits the full covariance without specifying the random-effects distribution and is expected to rank between ANCOVA and GLS.
\item A full comparison including both, and a feasible-GLS simulation under estimated components, is the natural continuation of this work.
\end{itemize}

\section{Data availability statement}\label{data-availability-statement}

\textbf{The package implementing all three estimators, the simulation drivers, and the shared bibliography are openly available in the project repository.}

\begin{itemize}\setlength{\itemsep}{2pt}\setlength{\parskip}{0pt}
\item The estimator implementations reside in the package source files.
\item The simulation drivers and shared compendium bibliography are included.
\item Specific paths appear in the Reproducibility section below.
\end{itemize}

\section{Reproducibility}\label{reproducibility}

\textbf{The manuscript was rendered with R 4.4.x using the package source and a small set of standard analysis and build packages.}

\begin{itemize}\setlength{\itemsep}{2pt}\setlength{\parskip}{0pt}
\item The in-package source supplies the estimator implementations.
\item The GLS comparator in the simulation is implemented directly in the manuscript source (the analyze functions of the simulation chunk); tinytest is the testing harness.
\item The manuscript is built with rmarkdown and bookdown.
\end{itemize}

\textbf{The simulation sets a single fixed seed before the replicate loop, which makes the run exactly reproducible.}

\begin{itemize}\setlength{\itemsep}{2pt}\setlength{\parskip}{0pt}
\item A single seed is set once at the entry point of the Monte Carlo comparison.
\item No per-replicate reseeding is performed; the replicate stream follows from the single seed.
\item Re-running the chunk therefore reproduces the tabulated results exactly.
\end{itemize}

\textbf{All code, drivers, companion manuscripts, and the shared bibliography are located at documented paths within the compendium.}

\begin{itemize}\setlength{\itemsep}{2pt}\setlength{\parskip}{0pt}
\item The package source and simulation drivers reside in their respective directories.
\item The three companion manuscripts are stored alongside this paper in the compendium.
\item The manuscript source and shared bibliography are co-located in the paper directory.
\end{itemize}

\end{document}
